\section{Resultados}

\subsection{Calibração da sonda de Hall}

Primeiramente, foi efetuada a calibração da sonda de efeito de Hall. Para tal, realizou-se uma montagem constítuida por uma fonte de $15 \, \unit{\volt}$, um reóstato, e um amperímetro. A sonda foi introduzida num solenóide-padrão, e posicionada no centro geométrico do mesmo para minimizar o erro associado à aplicação da equação \cref{eq:solenoid}, válida para solenóides para os quais o comprimento é muito maior que o raio e $N/L$ é conhecido:

\begin{equation}
	B = \mu_0 \, \frac{N}{L}I_S
	\label{eq:solenoid}
\end{equation}

em que $B$ corresponde ao campo magnético, $\mu_0$ à permeabilidade do vácuo, $N$ ao número de espiras, $L$ ao comprimento da solenóide e $I_S$ à corrente na solenóide.

Fez-se então variar a corrente que percorria o solenóide, $I_S$ através do reóstato, e mediu-se a diferença de potencial elétrico $V_H$ gerada nos terminais da sonda devido ao efeito desse campo, para cada valor de corrente. As medições  realizadas encontram-se na \cref{tab:calib-sonda}:


\begin{table}[!ht]
	\centering
	\caption{Tensão de Hall e corrente na solenóide medidas.}
	\label{tab:calib-sonda}
	% Use pgfplotstable to render the CSV
	\pgfplotstabletypeset[
		col sep=comma, % Tells LaTeX the file is a CSV
		% string type,   % Treat headers as text
		% column type={|c|},
		columns={Vh, Is},
		every head row/.style={
				before row=\hline,
				after row=\hline
			},
		every last row/.style={
				after row=\hline
			},
		% Assigning column headers and units
		columns/Vh/.style={
				column name={$V_H$ [\unit{mV}]},
				% column type={|c|},
				% column type={S[table-format=2.1]|} % Adjust format to your max digits
			},
		columns/Is/.style={
				column name={$I_S$ [\unit{mA}]},
				% column type={c|},
				% column type={|S[table-format=2.2]|} % Adjust format to your max digits
			},
	]{dados/solenoide.csv}
\end{table}


\subsection{Campo magnético em bobinas de Helmholtz}
Na segunda parte da atividade experimental, recorreu-se à sonda para medir o campo magnético e a sua variação face a alterações no circuito. Procurou-se também verificar o princípio da sobreposição de campos magnéticos.


Primeiramente, foi fechado um circuito entre a sonda e o multímetro com o intuito de verificar que a sonda não estava a sofrer interferência de tensões residuais. Para tal foi ajustado o potenciómetro da sonda até que a tensão indicada no multímetro fosse nula. De seguida, partindo da montagem anterior, substituiu-se o solenóide por duas bobinas de Helmholtz ligadas em série. Estas foram posicionadas a $6,475 \, \unit{\cm}$ uma da outra, o equivalente ao seu raio $R$, calculado como média aritmética do seu raio exterior e interior.
A sonda foi colocada de forma a atravessar o centro geométrico das bobinas. Fixou-se um valor de corrente de $I_S = 0,5 \, \unit{\ampere}$ para o resto da experiência e tendo em conta posição da escala da sonda relativamente à da bobina, registou-se uma posição de referência para sonda $x = 0$, e mediu-se a variação da tensão $V_H$ à medida que se distanciava  a sonda da sua posição inicial. Analogamente ao método adotado para a calibração da sonda, obtêm-se valores de $B$ para cada posição considerada.


Repetiu-se o procedimento para avaliar o efeito de cada bobina individualmente, e para uma outra configuração em que a corrente circulava em sentidos contrários nas duas bobinas. Na \crefrange{tab:Vh-serie}{tab:Vh-opostas} constam as medições efetuadas para as várias configurações consideradas.

\begin{table}[!ht]
	\centering
	\begin{minipage}{0.35\linewidth}
		\centering
		\caption{Bobinas em série}
		\label{tab:Vh-serie}
		\pgfplotstabletypeset[
			col sep=comma, % Tells LaTeX the file is a CSV
			% string type,   % Treat headers as text
			% column type={|c|},
			columns={x, Vh},
			every head row/.style={
					before row=\hline,
					after row=\hline
				},
			every last row/.style={
					after row=\hline
				},
			% Assigning column headers and units
			columns/x/.style={
					column name={$x$ [\unit{cm}]},
					% column type={c|},
					% column type={|S[table-format=2.2]|} % Adjust format to your max digits
				},
			columns/Vh/.style={
					column name={$V_H$ [\unit{mV}]},
					% column type={|c|},
					% column type={S[table-format=2.1]|} % Adjust format to your max digits
				},
		]{dados/bobinas-serie.csv}
	\end{minipage}
	\begin{minipage}{0.35\linewidth}
		\centering
		\caption{Bobina da direita}
		\label{tab:Vh-direita}
		\pgfplotstabletypeset[
			col sep=comma, % Tells LaTeX the file is a CSV
			% string type,   % Treat headers as text
			% column type={|c|},
			columns={x, Vh},
			every head row/.style={
					before row=\hline,
					after row=\hline
				},
			every last row/.style={
					after row=\hline
				},
			% Assigning column headers and units
			columns/x/.style={
					column name={$x$ [\unit{cm}]},
					% column type={c|},
					% column type={|S[table-format=2.2]|} % Adjust format to your max digits
				},
			columns/Vh/.style={
					column name={$V_H$ [\unit{mV}]},
					% column type={|c|},
					% column type={S[table-format=2.1]|} % Adjust format to your max digits
				},
		]{dados/bobins-rights.csv}
	\end{minipage}
\end{table}
\begin{table}[!ht]
	\centering
	\begin{minipage}{0.35\linewidth}
		\centering
		\caption{Bobina da esquerda}
		\label{tab:Vh-esquerda}
		\pgfplotstabletypeset[
			col sep=comma, % Tells LaTeX the file is a CSV
			% string type,   % Treat headers as text
			% column type={|c|},
			columns={x, Vh},
			every head row/.style={
					before row=\hline,
					after row=\hline
				},
			every last row/.style={
					after row=\hline
				},
			% Assigning column headers and units
			columns/x/.style={
					column name={$x$ [\unit{cm}]},
					% column type={c|},
					% column type={|S[table-format=2.2]|} % Adjust format to your max digits
				},
			columns/Vh/.style={
					column name={$V_H$ [\unit{mV}]},
					% column type={|c|},
					% column type={S[table-format=2.1]|} % Adjust format to your max digits
				},
		]{dados/lefts-bobains.csv}
	\end{minipage}
	\begin{minipage}{0.35\linewidth}
		\centering
		\caption{Bobinas em oposição de corrente}
		\label{tab:Vh-opostas}
		\pgfplotstabletypeset[
			col sep=comma, % Tells LaTeX the file is a CSV
			% string type,   % Treat headers as text
			% column type={|c|},
			columns={x, Vh},
			every head row/.style={
					before row=\hline,
					after row=\hline
				},
			every last row/.style={
					after row=\hline
				},
			% Assigning column headers and units
			columns/x/.style={
					column name={$x$ [\unit{cm}]},
					% column type={c|},
					% column type={|S[table-format=2.2]|} % Adjust format to your max digits
				},
			columns/Vh/.style={
					column name={$V_H$ [\unit{mV}]},
					% column type={|c|},
					% column type={S[table-format=2.1]|} % Adjust format to your max digits
				},
		]{dados/oposed-bobains.csv}
	\end{minipage}
\end{table}
